% Complete chapter-2 typesetting source; compact textbook presentation.
% Authoritative Markdown: ../../chapters/ch02/answer-key.md
% Stable task/source IDs are comments. All prose, tasks, tables, and notes retained.
% Local editorial-file links retain their labels; external links remain live.
% Shared layout for the two complete chapter-2 PDFs.
% This is a small hand-maintained TeX layout, not a conversion pipeline.
\RequirePackage{fix-cm}
\documentclass[11pt,twoside]{article}
\usepackage{fontspec}
\setmainfont{Cambria}[
  Renderer=HarfBuzz,
  BoldFont=cambriab.ttf,
  ItalicFont=cambriai.ttf,
  BoldItalicFont=cambriaz.ttf
]
\setmonofont{lmmono10-regular.otf}[BoldFont=lmmonolt10-bold.otf]
\usepackage[paperwidth=165mm,paperheight=241mm,inner=21mm,outer=19mm,
  top=22mm,bottom=20mm,headheight=14pt,headsep=7mm,footskip=10mm]{geometry}
\usepackage{xcolor,array,booktabs,longtable,tabularx}
\usepackage{enumitem,fancyhdr}
\usepackage[unicode,hidelinks]{hyperref}
\definecolor{ink}{HTML}{000000}
\definecolor{muted}{HTML}{4C5558}
\setlength{\parindent}{0pt}
\setlength{\parskip}{3.5pt}
\AtBeginDocument{\fontsize{11}{14.3}\selectfont}
\setlength{\tabcolsep}{4pt}
\renewcommand{\arraystretch}{1.10}
\setlist[itemize]{leftmargin=1.3em,label=\textbullet,itemsep=2pt,topsep=3pt,parsep=0pt}
\setlist[enumerate]{leftmargin=1.8em,itemsep=3pt,topsep=4pt,parsep=0pt}
\widowpenalty=10000
\clubpenalty=10000
\emergencystretch=1em
\raggedbottom
\pagestyle{fancy}
\fancyhf{}
\fancyhead[LE,RO]{\small\itshape Plural nouns}
\fancyhead[LO,RE]{\small Cook Expanded}
\fancyfoot[LE,RO]{\small\thepage}
\renewcommand{\headrulewidth}{0.3pt}
\newcommand{\booktitle}[2]{
  {\fontsize{9}{11}\selectfont\color{muted}#1\par}
  \vspace{3pt}
  {\fontsize{20}{24}\selectfont\color{ink}#2\par}
  \vspace{6pt}
}
\newcommand{\source}[1]{\par{\footnotesize\color{muted}#1\par}}
\newcommand{\readingtext}[1]{
  \par\addvspace{5pt}
  \begingroup\leftskip=1.5em
    \fontsize{12}{16}\selectfont #1\par
  \endgroup\addvspace{5pt}
}
\makeatletter
\renewcommand{\section}{\@startsection{section}{1}{\z@}
  {-2.4ex \@plus -0.6ex \@minus -0.2ex}{0.8ex \@plus 0.2ex}
  {\normalfont\fontsize{12}{15}\selectfont\bfseries\color{ink}}}
\renewcommand{\subsection}{\@startsection{subsection}{2}{\z@}
  {-1.8ex \@plus -0.5ex \@minus -0.2ex}{0.5ex \@plus 0.2ex}
  {\normalfont\fontsize{11}{14}\selectfont\bfseries\color{ink}}}
\makeatother

\hypersetup{pdftitle={Chapter 2: Answers and commentary},pdfauthor={Cook Expanded}}
\fancyfoot[LO,RE]{\footnotesize\color{muted}Answers and commentary}
\begin{document}

\booktitle{CHAPTER 2 · ANSWER KEY}{Answers and commentary}

% Chapter: ch02. Independently reviewed. Stable item IDs match coverage.json.

\pdfbookmark[1]{Preparation (§3)}{key-section-1}
\section*{Preparation (§3)}


\begin{enumerate}
% item: ch02-prep-01
\item[1.] \textbf{ðā dagas}; \textbf{ðā wyrta}; \textbf{ðā dēor}: “the days”; “the herbs”; “the animals.” Each phrase is nominative plural. The article becomes \textbf{ðā} in all three; \textbf{dēor} itself is unchanged.
% item: ch02-prep-02
\item[2.] \textbf{fuglas}, “birds”; \textbf{tācnu}, “signs”; \textbf{męnn}, “men.” The first two lose an unstressed stem vowel; the last changes its root vowel.
% item: ch02-prep-03
\item[3.] \textbf{sǣ}, “sea,” masculine (also feminine); \textbf{hwæl}, “whale,” masculine; \textbf{nīeten}, “animal; beast,” neuter, with plural \textbf{nīetenu} used for “cattle.” The task tests recognition of the headwords and their genders.
\end{enumerate}


\pdfbookmark[1]{Composition (§4)}{key-section-2}
\section*{Composition (§4)}


The target is two definite nominative plural noun phrases joined by \textbf{and}. Each takes \textbf{ðā}, regardless of gender. The noun forms follow the selected paradigms in §§1–2.


\begin{longtable}{@{}>{\raggedright\arraybackslash}p{8mm}>{\raggedright\arraybackslash}p{63mm}>{\raggedright\arraybackslash}p{48mm}@{}}
\toprule
\textbf{No.} & \textbf{Old English} & \textbf{Meaning} \\
\midrule
\endfirsthead
\toprule
\textbf{No.} & \textbf{Old English} & \textbf{Meaning} \\
\midrule
\endhead
\midrule
\multicolumn{3}{r}{\footnotesize\itshape Continued on the next page} \\
\endfoot
\bottomrule
\endlastfoot
% item: EX-I-02-01
1 & \textbf{ðā dagas and ðā niht} & The days and the nights. \\
% item: EX-I-02-02
2 & \textbf{ðā fiscas and ðā fuglas} & The fishes and the birds. \\
% item: EX-I-02-03
3 & \textbf{ðā trēowu and ðā wyrta} & The trees and the herbs. \\
% item: EX-I-02-04
4 & \textbf{ðā andlīcnessa and ðā gelīcnessa} & The images and the likenesses. \\
% item: EX-I-02-05
5 & \textbf{ðā tācnu and ðā tīda} & The signs and the seasons. \\
% item: EX-I-02-06
6 & \textbf{ðā gēar and ðā dagas} & The years and the days. \\
% item: EX-I-02-07
7 & \textbf{ðā męnn and ðā dēor} & The men and the animals. \\
% item: EX-I-02-08
8 & \textbf{ðā sǣs and ðā heofonas} & The seas and the heavens. \\
% item: EX-I-02-09
9 & \textbf{ðā sǣd and ðā wæstmas} & The seeds and the fruits. \\
% item: EX-I-02-10
10 & \textbf{ðā ǣfenu and ðā morgnas} & The evenings and the mornings. \\
% item: EX-I-02-11
11 & \textbf{ðā gāstas and ðā męnn} & The spirits and the men. \\
% item: EX-I-02-12
12 & \textbf{ðā hwalas and ðā nīetenu} & The whales and the cattle. \\
\end{longtable}


\textbf{Variants.} In 2, \textbf{fugolas}; in 3–5, feminine \textbf{wyrte}, \textbf{andlīcnesse}, \textbf{gelīcnesse}, \textbf{tīde} beside the selected \textbf{-a} forms; in 4, \textbf{anlīcnessa}; in 8, weak feminine \textbf{heofonan}; in 10, \textbf{morgenas}. These are grammatical alternatives with the required meanings, though they differ from the selected vocabulary forms. The alternative \textbf{heofonan} belongs to singular \textbf{heofone}, not the selected masculine paradigm. In 7, taught \textbf{nīetenu} is also acceptable for “animals.”


\textbf{Form and meaning.} \textbf{niht}, \textbf{gēar}, \textbf{dēor}, and \textbf{sǣd} in these answers are plural despite matching the singular spelling. \textbf{dēor} means animals generally; \textbf{nīetenu} here means cattle. \textbf{ǣfenu} retains the selected neuter gender. The \textbf{æ/a} alternation in \textbf{dagas, hwalas} does not apply to \textbf{wæstmas}, which retains short \textbf{æ}.


\pdfbookmark[1]{Reading (§5)}{key-section-3}
\section*{Reading (§5)}


\begin{enumerate}
% item: ch02-reading-01
\item[1.] The birds are to increase in number over or across the earth. This is the content of a command or wish in the blessing.
% item: ch02-reading-02
\item[2.] \textbf{ðā} (article) + \textbf{fuglas} (noun): masculine nominative plural, the subject of \textbf{bēon gemanigfielde}. The article has the same plural form for every gender; the noun is masculine.
% item: ch02-reading-03
\item[3.] Close: “and let the birds be multiplied over the earth.” Fluent: “and let the birds multiply across the earth.” “May the birds increase in number throughout the earth” also preserves the optative sense; a close translation retains \textbf{and}.
\end{enumerate}


English intransitive “multiply” conveys the meaning of the Old English passive. “The birds multiplied” or a simple prediction “the birds will multiply” loses the command or wish. \textbf{bēon} is a finite plural optative here, not an infinitive. No unique accusative/dative analysis of \textbf{eorðan} is required; the entire phrase is supplied.


\pdfbookmark[1]{Review (§6)}{key-section-4}
\section*{Review (§6)}


\begin{enumerate}
% item: ch02-review-01
\item[1.] \textbf{ðā sǣd and ðā fuglas}, “the seeds and the birds”: neuter and masculine nominative plural. \textbf{fugolas} is a valid alternative.
% item: ch02-review-02
\item[2.] \textbf{ðā hwalas and ðā sǣs}, “the whales and the seas”: nominative plural; both nouns use the selected masculine gender. \textbf{hwalas} has the stem vowel \textbf{a}; \textbf{sǣs} is the individually taught plural.
\begin{samepage}
% item: ch02-review-03
\item[3.] \textbf{ðæt sǣd}, neuter singular; \textbf{ðā sǣd}, neuter plural. The article distinguishes number. \textbf{sǣd} has long \textbf{ǣ}: as a monosyllabic strong neuter with a heavy stem, it regularly has no nominative plural ending.
\par
\end{samepage}
\end{enumerate}


\pdfbookmark[1]{Assessment notes}{key-section-5}
\section*{Assessment notes}


Separate grammatical validity, meaning, and compliance with the specified vocabulary and construction. Verified variants need not match the model string. Equivalent explanations and standard grammatical abbreviations are sufficient for analytical answers.


Common errors are retaining a singular article before a plural noun, adding \textbf{-as} to every noun, overlooking the stem changes, and treating an unchanged noun as necessarily singular. The table supplies the required forms; gender alone does not predict them.


Capitalization and final punctuation do not determine correctness. Thorn/\allowbreak{}eth alternatives are acceptable. Missing length marks and substitution of ordinary \textbf{e} for the retained \textbf{ę} should be corrected as spelling or editorial-convention differences, separately from grammar.


\begingroup\small
\pdfbookmark[1]{Notes on sources and variants}{key-section-6}
\section*{Notes on sources and variants}


\textbf{Ex:} A. S. Cook, \emph{Exercises in Old English} (1895). \textbf{FB:} \emph{A First Book in Old English}, 2nd ed. (1895). References use printed pages. Answers and translations are newly prepared; full evidence is in the editorial record.


\begin{enumerate}
\item[1.] \textbf{Production and review tasks:} Ex I.2, p.1, and vocabulary; FB §§4, 17, 43, 46–47, 51–52, 84; \textbf{nīeten}, FB vocabulary p.307. Stem weight and neuter endings: §4, p.6; §47, pp.32–33. \textbf{ǣfenu} is formed from the neuter patterns in §47.1, 7, not quoted from a located passage. \textbf{tācnu}: \href{https://bosworthtoller.com/58972}{Bosworth–Toller}. \textbf{sǣs}: FB p.124, note 25, and \href{https://bosworthtoller.com/26140}{Bosworth–Toller}.
\item[2.] \textbf{Variants:} \textbf{fugolas}, \href{https://bosworthtoller.com/12620}{Bosworth–Toller, \emph{fugol}}; \textbf{morgenas}, \href{https://bosworthtoller.com/23129}{\emph{morgen}}; feminine \textbf{-a/-e}, FB §51, pp.35–36; \textbf{an(d)līcnes}, Ex p.51; \textbf{heofone/heofonan}, FB §53, p.37.
\item[3.] \textbf{Correction C01:} \textbf{wæstm, wæstmas} retain short \textbf{æ}: \href{https://glossary.oldenglishaerobics.net/oeaginfo.php?objid=xwaestm\&objtype=gentry}{Baker's glossary}; \href{https://doi.org/10.17811/selim.29.2024.109-25}{Stenroos (2024), p.118, Table 1}.
\item[4.] \textbf{Reading:} FB I, p.126, printed lines 12–13, notes 19 and 21; §138, p.78; §193(a), p.103; vocabulary p.293. Weak \textbf{eorðe/eorðan}: §53, p.37; spatial \textbf{ofer} with accusative or dative: \href{https://bosworthtoller.com/24000}{Bosworth–Toller}.
\end{enumerate}

\endgroup

\end{document}
